\documentclass[12pt,a4paper]{article}

\usepackage[utf8]{inputenc}
\usepackage[T1]{fontenc}
\usepackage{amsmath, amssymb}
\usepackage{enumitem}
\usepackage{geometry}
\usepackage{xcolor}
\usepackage{multicol}

\geometry{margin=2.3cm}

\title{Banco de Questões Objetivas -- Aula 10\\SQL}
\author{Banco de Dados}
\date{}

\begin{document}

\maketitle

\section*{Instruções}

Este banco contém questões objetivas sobre os principais tópicos da Aula 10:
\begin{itemize}
    \item SQL como linguagem declarativa;
    \item SQL padrão e portabilidade entre SGBDs;
    \item DDL, DML e VDL;
    \item criação de esquemas e tabelas;
    \item tipos de dados;
    \item restrições de integridade;
    \item consultas com \texttt{SELECT}, \texttt{FROM} e \texttt{WHERE};
    \item uso de \texttt{DISTINCT};
    \item junções;
    \item consultas aninhadas;
    \item funções agregadas;
    \item \texttt{GROUP BY} e \texttt{HAVING};
    \item comandos \texttt{INSERT}, \texttt{DELETE} e \texttt{UPDATE}.
\end{itemize}

\section{Questões de Múltipla Escolha}

\begin{enumerate}[label=\textbf{Q\arabic*.}]

\item SQL significa:

\begin{enumerate}[label=\alph*)]
    \item Structured Query Language.
    \item Simple Query Logic.
    \item System Query Level.
    \item Sequential Query Language.
\end{enumerate}

\item SQL é considerada uma linguagem declarativa porque:

\begin{enumerate}[label=\alph*)]
    \item o usuário especifica passo a passo como o SGBD deve executar a consulta.
    \item o usuário especifica o resultado desejado, e o SGBD decide como executar.
    \item ela só permite comandos de criação de tabela.
    \item ela não permite restrições de integridade.
\end{enumerate}

\item Uma vantagem de manter-se próximo ao SQL padrão é:

\begin{enumerate}[label=\alph*)]
    \item impedir a migração entre SGBDs.
    \item facilitar migrações entre SGBDs diferentes.
    \item eliminar a necessidade de chaves primárias.
    \item proibir consultas com \texttt{JOIN}.
\end{enumerate}

\item A sigla DDL refere-se a comandos de:

\begin{enumerate}[label=\alph*)]
    \item definição de dados.
    \item manipulação de dados.
    \item definição de visões apenas.
    \item controle de concorrência apenas.
\end{enumerate}

\item A sigla DML refere-se principalmente a comandos de:

\begin{enumerate}[label=\alph*)]
    \item definição de tabelas e esquemas.
    \item consulta, inserção, remoção e atualização de dados.
    \item criação de usuários do sistema operacional.
    \item definição exclusiva de visões.
\end{enumerate}

\item A sigla VDL está relacionada à:

\begin{enumerate}[label=\alph*)]
    \item definição de visões.
    \item remoção de bancos de dados.
    \item atualização de tuplas.
    \item criação de índices físicos obrigatórios.
\end{enumerate}

\item No vocabulário de SQL, uma relação do modelo relacional corresponde a:

\begin{enumerate}[label=\alph*)]
    \item uma tabela.
    \item uma tupla.
    \item uma coluna.
    \item uma restrição.
\end{enumerate}

\item No vocabulário de SQL, uma tupla corresponde a:

\begin{enumerate}[label=\alph*)]
    \item uma tabela.
    \item uma linha.
    \item uma coluna.
    \item um domínio.
\end{enumerate}

\item No vocabulário de SQL, um atributo corresponde a:

\begin{enumerate}[label=\alph*)]
    \item uma tabela.
    \item uma linha.
    \item uma coluna.
    \item um esquema inteiro.
\end{enumerate}

\item O comando \texttt{CREATE} pode ser usado para criar:

\begin{enumerate}[label=\alph*)]
    \item apenas linhas em uma tabela.
    \item esquemas, tabelas, domínios, visões, asserções e triggers.
    \item apenas consultas temporárias.
    \item apenas usuários do sistema operacional.
\end{enumerate}

\item Um esquema SQL é:

\begin{enumerate}[label=\alph*)]
    \item um agrupamento de elementos como tabelas, restrições, visões e domínios.
    \item uma linha específica de uma tabela.
    \item sempre uma chave primária.
    \item um comando usado apenas para apagar dados.
\end{enumerate}

\item A estrutura básica de uma consulta SQL é:

\begin{enumerate}[label=\alph*)]
    \item \texttt{SELECT-FROM-WHERE}.
    \item \texttt{CREATE-FROM-WHERE}.
    \item \texttt{INSERT-GROUP-HAVING}.
    \item \texttt{DROP-SELECT-TABLE}.
\end{enumerate}

\item Na consulta SQL, a cláusula \texttt{SELECT} indica:

\begin{enumerate}[label=\alph*)]
    \item as tabelas utilizadas na consulta.
    \item os atributos ou expressões que aparecerão no resultado.
    \item a condição de filtragem das linhas.
    \item a ordem física das tabelas no disco.
\end{enumerate}

\item Na consulta SQL, a cláusula \texttt{FROM} indica:

\begin{enumerate}[label=\alph*)]
    \item as tabelas envolvidas na consulta.
    \item os atributos projetados no resultado.
    \item os valores a serem atualizados.
    \item os grupos a serem removidos.
\end{enumerate}

\item Na consulta SQL, a cláusula \texttt{WHERE} indica:

\begin{enumerate}[label=\alph*)]
    \item a condição usada para selecionar as tuplas.
    \item a criação de uma nova tabela.
    \item a definição de uma chave primária.
    \item a ordem final de apresentação.
\end{enumerate}

\item Se uma consulta possui várias tabelas no \texttt{FROM} e não possui condição de junção no \texttt{WHERE}, o resultado tende a ser:

\begin{enumerate}[label=\alph*)]
    \item uma chave primária.
    \item um produto cartesiano.
    \item uma visão materializada.
    \item uma restrição de domínio.
\end{enumerate}

\item O operador \texttt{*} em uma consulta como:

\[
\texttt{SELECT * FROM Aluno;}
\]

significa:

\begin{enumerate}[label=\alph*)]
    \item retornar todos os atributos das tuplas selecionadas.
    \item retornar apenas a chave primária.
    \item remover todas as tuplas.
    \item criar uma nova tabela.
\end{enumerate}

\item Em SQL, por padrão, o resultado de uma consulta pode conter tuplas duplicadas. Para removê-las, usa-se:

\begin{enumerate}[label=\alph*)]
    \item \texttt{UNIQUE TABLE}.
    \item \texttt{DISTINCT}.
    \item \texttt{PRIMARY}.
    \item \texttt{REMOVE DUPLICATES}.
\end{enumerate}

\item A consulta:

\[
\texttt{SELECT DISTINCT curso FROM Aluno;}
\]

retorna:

\begin{enumerate}[label=\alph*)]
    \item todos os cursos, incluindo repetições.
    \item apenas cursos distintos.
    \item todas as colunas da tabela Aluno.
    \item apenas alunos sem curso.
\end{enumerate}

\item A cláusula \texttt{ORDER BY} é usada para:

\begin{enumerate}[label=\alph*)]
    \item ordenar o resultado da consulta.
    \item agrupar tuplas.
    \item criar restrições de integridade.
    \item eliminar uma tabela do banco.
\end{enumerate}

\item A cláusula \texttt{GROUP BY} é usada para:

\begin{enumerate}[label=\alph*)]
    \item agrupar tuplas com base em um ou mais atributos.
    \item remover todas as linhas duplicadas obrigatoriamente.
    \item criar uma chave estrangeira.
    \item atualizar várias tabelas ao mesmo tempo.
\end{enumerate}

\item A cláusula \texttt{HAVING} é usada para:

\begin{enumerate}[label=\alph*)]
    \item filtrar grupos após o agrupamento.
    \item filtrar linhas antes do agrupamento.
    \item criar uma tabela.
    \item definir um tipo de dado.
\end{enumerate}

\item A diferença principal entre \texttt{WHERE} e \texttt{HAVING} é:

\begin{enumerate}[label=\alph*)]
    \item \texttt{WHERE} filtra tuplas; \texttt{HAVING} filtra grupos.
    \item \texttt{WHERE} filtra grupos; \texttt{HAVING} cria tabelas.
    \item ambas fazem exatamente a mesma coisa em qualquer contexto.
    \item \texttt{HAVING} só pode ser usado com \texttt{INSERT}.
\end{enumerate}

\item Qual das opções abaixo contém apenas funções agregadas?

\begin{enumerate}[label=\alph*)]
    \item \texttt{COUNT}, \texttt{SUM}, \texttt{MAX}, \texttt{MIN}, \texttt{AVG}.
    \item \texttt{SELECT}, \texttt{FROM}, \texttt{WHERE}.
    \item \texttt{CREATE}, \texttt{DROP}, \texttt{ALTER}.
    \item \texttt{INSERT}, \texttt{DELETE}, \texttt{UPDATE}.
\end{enumerate}

\item A função \texttt{COUNT(*)} retorna:

\begin{enumerate}[label=\alph*)]
    \item a soma dos valores de uma coluna.
    \item o número de linhas selecionadas.
    \item o maior valor de uma coluna.
    \item a média dos valores de uma coluna.
\end{enumerate}

\item A função \texttt{AVG(salario)} retorna:

\begin{enumerate}[label=\alph*)]
    \item o maior salário.
    \item o menor salário.
    \item a média dos salários.
    \item a quantidade de salários distintos.
\end{enumerate}

\item A função \texttt{SUM(salario)} retorna:

\begin{enumerate}[label=\alph*)]
    \item a soma dos salários.
    \item a média dos salários.
    \item o menor salário.
    \item o número de tabelas do banco.
\end{enumerate}

\item A função \texttt{MAX(salario)} retorna:

\begin{enumerate}[label=\alph*)]
    \item o maior salário.
    \item o menor salário.
    \item a média dos salários.
    \item a quantidade de salários.
\end{enumerate}

\item A função \texttt{MIN(salario)} retorna:

\begin{enumerate}[label=\alph*)]
    \item o maior salário.
    \item o menor salário.
    \item a média dos salários.
    \item a soma dos salários.
\end{enumerate}

\item Considere a consulta:

\[
\texttt{SELECT curso, COUNT(*) FROM Aluno GROUP BY curso;}
\]

Ela retorna:

\begin{enumerate}[label=\alph*)]
    \item a quantidade de alunos por curso.
    \item todos os alunos ordenados por matrícula.
    \item todos os cursos sem contagem.
    \item apenas alunos com curso nulo.
\end{enumerate}

\item Considere a consulta:

\[
\texttt{SELECT curso, COUNT(*) FROM Aluno GROUP BY curso HAVING COUNT(*) > 10;}
\]

Ela retorna:

\begin{enumerate}[label=\alph*)]
    \item todos os alunos com mais de 10 anos.
    \item cursos com mais de 10 alunos.
    \item alunos cujo nome tem mais de 10 letras.
    \item todas as tabelas com mais de 10 colunas.
\end{enumerate}

\item O comando \texttt{INSERT} é usado para:

\begin{enumerate}[label=\alph*)]
    \item inserir novas tuplas em uma tabela.
    \item remover uma tabela inteira do catálogo.
    \item alterar a estrutura de um esquema.
    \item definir uma chave estrangeira.
\end{enumerate}

\item O comando \texttt{DELETE} é usado para:

\begin{enumerate}[label=\alph*)]
    \item remover tuplas de uma tabela.
    \item criar um novo esquema.
    \item alterar o nome de uma coluna.
    \item calcular médias.
\end{enumerate}

\item O comando \texttt{UPDATE} é usado para:

\begin{enumerate}[label=\alph*)]
    \item modificar valores de atributos em tuplas existentes.
    \item criar uma nova tabela.
    \item apagar a definição da tabela.
    \item criar uma visão.
\end{enumerate}

\item O que acontece em um comando \texttt{DELETE} sem cláusula \texttt{WHERE}?

\begin{enumerate}[label=\alph*)]
    \item Todas as tuplas da tabela são removidas.
    \item A tabela é obrigatoriamente apagada do catálogo.
    \item Apenas a primeira tupla é removida.
    \item Nenhuma tupla pode ser removida.
\end{enumerate}

\item O que acontece em um comando \texttt{UPDATE} sem cláusula \texttt{WHERE}?

\begin{enumerate}[label=\alph*)]
    \item Todas as tuplas da tabela podem ser atualizadas.
    \item Nenhuma tupla é atualizada.
    \item A tabela é apagada.
    \item Uma nova tabela é criada.
\end{enumerate}

\item O comando \texttt{DROP TABLE} é usado para:

\begin{enumerate}[label=\alph*)]
    \item remover a definição da tabela do banco.
    \item remover apenas uma linha da tabela.
    \item atualizar todas as tuplas.
    \item fazer uma junção entre duas tabelas.
\end{enumerate}

\item O comando \texttt{ALTER TABLE} é usado para:

\begin{enumerate}[label=\alph*)]
    \item modificar a estrutura de uma tabela.
    \item consultar todos os dados de uma tabela.
    \item inserir uma nova linha.
    \item ordenar o resultado de uma consulta.
\end{enumerate}

\item Uma chave primária em SQL serve para:

\begin{enumerate}[label=\alph*)]
    \item identificar unicamente as tuplas de uma tabela.
    \item permitir valores duplicados obrigatoriamente.
    \item armazenar apenas valores nulos.
    \item impedir qualquer consulta na tabela.
\end{enumerate}

\item Uma chave estrangeira em SQL serve para:

\begin{enumerate}[label=\alph*)]
    \item representar uma referência a uma chave de outra tabela.
    \item identificar sempre uma coluna multivalorada.
    \item substituir a cláusula \texttt{WHERE}.
    \item ordenar automaticamente todas as consultas.
\end{enumerate}

\item A restrição \texttt{NOT NULL} indica que:

\begin{enumerate}[label=\alph*)]
    \item o atributo não pode receber valor nulo.
    \item o atributo deve ser sempre uma chave estrangeira.
    \item a tabela não pode ter linhas.
    \item o atributo não pode ser consultado.
\end{enumerate}

\item A restrição \texttt{UNIQUE} indica que:

\begin{enumerate}[label=\alph*)]
    \item os valores de um atributo ou conjunto de atributos não podem se repetir.
    \item todos os valores devem ser nulos.
    \item a tabela não pode ter chave primária.
    \item a consulta deve retornar todos os atributos.
\end{enumerate}

\item A restrição \texttt{CHECK} é usada para:

\begin{enumerate}[label=\alph*)]
    \item definir uma condição que os valores devem satisfazer.
    \item criar uma consulta aninhada obrigatória.
    \item remover tuplas duplicadas do resultado.
    \item apagar uma tabela do banco.
\end{enumerate}

\item Um \texttt{JOIN} é usado para:

\begin{enumerate}[label=\alph*)]
    \item combinar tuplas de tabelas relacionadas.
    \item apagar todos os dados de uma tabela.
    \item criar uma chave primária automaticamente.
    \item impedir o uso de agregação.
\end{enumerate}

\item Em uma junção entre \texttt{Aluno} e \texttt{Curso}, a condição:

\[
\texttt{Aluno.idCurso = Curso.idCurso}
\]

representa:

\begin{enumerate}[label=\alph*)]
    \item uma condição de junção.
    \item uma função agregada.
    \item uma cláusula de agrupamento.
    \item uma remoção de tabela.
\end{enumerate}

\item Uma consulta aninhada é:

\begin{enumerate}[label=\alph*)]
    \item uma consulta dentro de outra consulta.
    \item uma tabela sem chave primária.
    \item uma restrição de domínio.
    \item um comando exclusivo de remoção.
\end{enumerate}

\item O operador \texttt{IN} pode ser usado para:

\begin{enumerate}[label=\alph*)]
    \item verificar se um valor pertence a um conjunto de valores.
    \item apagar uma tabela.
    \item criar uma chave primária.
    \item ordenar resultados.
\end{enumerate}

\item Uma visão em SQL pode ser entendida como:

\begin{enumerate}[label=\alph*)]
    \item uma tabela virtual derivada de uma consulta.
    \item uma tabela física obrigatória com dados duplicados.
    \item uma chave estrangeira.
    \item uma restrição de domínio.
\end{enumerate}

\item Uma trigger é, em geral:

\begin{enumerate}[label=\alph*)]
    \item uma ação automática executada quando certos eventos ocorrem no banco.
    \item uma coluna que nunca aceita valores nulos.
    \item uma função agregada.
    \item uma forma de ordenar tuplas.
\end{enumerate}

\item Uma asserção em SQL é usada para:

\begin{enumerate}[label=\alph*)]
    \item especificar restrições gerais sobre o banco.
    \item inserir várias tuplas obrigatoriamente.
    \item criar uma tabela temporária sempre.
    \item executar uma consulta sem \texttt{FROM}.
\end{enumerate}

\item Em SQL, uma consulta pode conter, na forma mais completa apresentada, as cláusulas:

\begin{enumerate}[label=\alph*)]
    \item \texttt{SELECT}, \texttt{FROM}, \texttt{WHERE}, \texttt{GROUP BY}, \texttt{HAVING}, \texttt{ORDER BY}.
    \item \texttt{INSERT}, \texttt{DELETE}, \texttt{UPDATE}, \texttt{DROP}.
    \item \texttt{CREATE}, \texttt{TABLE}, \texttt{COLUMN}, \texttt{ROW}.
    \item \texttt{PRIMARY}, \texttt{FOREIGN}, \texttt{CHECK}, \texttt{UNIQUE}.
\end{enumerate}

\end{enumerate}

\section{Questões de Verdadeiro ou Falso}

\begin{enumerate}[resume,label=\textbf{Q\arabic*.}]

\item Julgue as afirmações:

\begin{enumerate}[label=\alph*)]
    \item ( \quad ) SQL é uma linguagem declarativa.
    \item ( \quad ) SQL foi projetada no contexto do modelo relacional.
    \item ( \quad ) SQL obriga o usuário a especificar o algoritmo físico de execução da consulta.
    \item ( \quad ) SQL tornou-se uma linguagem padrão para bancos relacionais.
\end{enumerate}

\item Julgue as afirmações:

\begin{enumerate}[label=\alph*)]
    \item ( \quad ) DDL está relacionada à definição de dados.
    \item ( \quad ) DML está relacionada a consultas e atualizações.
    \item ( \quad ) VDL está relacionada à definição de visões.
    \item ( \quad ) SQL não permite restrições de integridade.
\end{enumerate}

\item Julgue as afirmações:

\begin{enumerate}[label=\alph*)]
    \item ( \quad ) Em SQL, relação corresponde a tabela.
    \item ( \quad ) Em SQL, tupla corresponde a linha.
    \item ( \quad ) Em SQL, atributo corresponde a coluna.
    \item ( \quad ) Em SQL, tabela corresponde sempre a um banco inteiro.
\end{enumerate}

\item Julgue as afirmações:

\begin{enumerate}[label=\alph*)]
    \item ( \quad ) A cláusula \texttt{SELECT} indica os atributos ou expressões retornadas.
    \item ( \quad ) A cláusula \texttt{FROM} indica as tabelas envolvidas.
    \item ( \quad ) A cláusula \texttt{WHERE} indica condições de filtragem.
    \item ( \quad ) A cláusula \texttt{WHERE} é usada para filtrar grupos após agregação.
\end{enumerate}

\item Julgue as afirmações:

\begin{enumerate}[label=\alph*)]
    \item ( \quad ) \texttt{COUNT}, \texttt{SUM}, \texttt{MAX}, \texttt{MIN} e \texttt{AVG} são funções agregadas.
    \item ( \quad ) \texttt{GROUP BY} é usado para agrupar tuplas.
    \item ( \quad ) \texttt{HAVING} é usado para filtrar grupos.
    \item ( \quad ) \texttt{ORDER BY} é usado para criar tabelas.
\end{enumerate}

\item Julgue as afirmações:

\begin{enumerate}[label=\alph*)]
    \item ( \quad ) \texttt{INSERT} insere tuplas.
    \item ( \quad ) \texttt{DELETE} remove tuplas.
    \item ( \quad ) \texttt{UPDATE} altera valores de tuplas existentes.
    \item ( \quad ) \texttt{DROP TABLE} remove apenas uma tupla da tabela.
\end{enumerate}

\item Julgue as afirmações:

\begin{enumerate}[label=\alph*)]
    \item ( \quad ) Uma chave primária identifica unicamente tuplas.
    \item ( \quad ) Uma chave estrangeira representa referência entre tabelas.
    \item ( \quad ) \texttt{NOT NULL} permite obrigatoriamente valores nulos.
    \item ( \quad ) \texttt{CHECK} pode definir uma condição sobre valores permitidos.
\end{enumerate}

\item Julgue as afirmações:

\begin{enumerate}[label=\alph*)]
    \item ( \quad ) Um \texttt{JOIN} combina dados de tabelas relacionadas.
    \item ( \quad ) Uma consulta aninhada é uma consulta dentro de outra.
    \item ( \quad ) Uma visão pode ser entendida como uma tabela virtual.
    \item ( \quad ) Uma trigger é uma função agregada.
\end{enumerate}

\end{enumerate}

\section{Questões de Aplicação Objetiva}

\begin{enumerate}[resume,label=\textbf{Q\arabic*.}]

\item Considere a tabela:

\[
Aluno(\text{matricula}, \text{nome}, \text{curso})
\]

Qual consulta retorna o nome de todos os alunos do curso \texttt{'BTI'}?

\begin{enumerate}[label=\alph*)]
    \item \texttt{SELECT nome FROM Aluno WHERE curso = 'BTI';}
    \item \texttt{FROM nome SELECT Aluno WHERE curso = 'BTI';}
    \item \texttt{SELECT curso FROM nome WHERE Aluno = 'BTI';}
    \item \texttt{DELETE nome FROM Aluno WHERE curso = 'BTI';}
\end{enumerate}

\item Considere:

\[
Aluno(\text{matricula}, \text{nome}, \text{idCurso})
\]
\[
Curso(\text{idCurso}, \text{nomeCurso})
\]

Qual consulta retorna o nome do aluno e o nome do curso correspondente?

\begin{enumerate}[label=\alph*)]
    \item \texttt{SELECT Aluno.nome, Curso.nomeCurso FROM Aluno, Curso WHERE Aluno.idCurso = Curso.idCurso;}
    \item \texttt{SELECT Aluno.nome, Curso.nomeCurso FROM Aluno, Curso;}
    \item \texttt{DELETE FROM Aluno WHERE idCurso = Curso.idCurso;}
    \item \texttt{UPDATE Aluno SET nomeCurso = Curso.nomeCurso;}
\end{enumerate}

\item Qual consulta retorna a quantidade de alunos por curso?

\begin{enumerate}[label=\alph*)]
    \item \texttt{SELECT curso, COUNT(*) FROM Aluno GROUP BY curso;}
    \item \texttt{SELECT curso FROM Aluno WHERE COUNT(*) GROUP BY curso;}
    \item \texttt{SELECT COUNT(curso) FROM Aluno WHERE GROUP BY curso;}
    \item \texttt{DELETE curso, COUNT(*) FROM Aluno;}
\end{enumerate}

\item Qual consulta retorna apenas cursos com mais de 5 alunos?

\begin{enumerate}[label=\alph*)]
    \item \texttt{SELECT curso, COUNT(*) FROM Aluno GROUP BY curso HAVING COUNT(*) > 5;}
    \item \texttt{SELECT curso, COUNT(*) FROM Aluno WHERE COUNT(*) > 5;}
    \item \texttt{SELECT curso FROM Aluno HAVING curso > 5;}
    \item \texttt{SELECT curso FROM Aluno GROUP BY COUNT(*) > 5;}
\end{enumerate}

\item Qual comando insere um aluno na tabela abaixo?

\[
Aluno(\text{matricula}, \text{nome}, \text{curso})
\]

\begin{enumerate}[label=\alph*)]
    \item \texttt{INSERT INTO Aluno VALUES (1, 'Ana', 'BTI');}
    \item \texttt{DELETE INTO Aluno VALUES (1, 'Ana', 'BTI');}
    \item \texttt{UPDATE INTO Aluno VALUES (1, 'Ana', 'BTI');}
    \item \texttt{CREATE INTO Aluno VALUES (1, 'Ana', 'BTI');}
\end{enumerate}

\item Qual comando atualiza o curso do aluno de matrícula 10 para \texttt{'Computacao'}?

\begin{enumerate}[label=\alph*)]
    \item \texttt{UPDATE Aluno SET curso = 'Computacao' WHERE matricula = 10;}
    \item \texttt{UPDATE Aluno WHERE curso = 'Computacao' SET matricula = 10;}
    \item \texttt{DELETE Aluno SET curso = 'Computacao' WHERE matricula = 10;}
    \item \texttt{INSERT Aluno SET curso = 'Computacao' WHERE matricula = 10;}
\end{enumerate}

\item Qual comando remove o aluno de matrícula 10?

\begin{enumerate}[label=\alph*)]
    \item \texttt{DELETE FROM Aluno WHERE matricula = 10;}
    \item \texttt{DROP FROM Aluno WHERE matricula = 10;}
    \item \texttt{UPDATE FROM Aluno WHERE matricula = 10;}
    \item \texttt{REMOVE TABLE Aluno WHERE matricula = 10;}
\end{enumerate}

\item Qual comando cria uma tabela simples de aluno com matrícula como chave primária?

\begin{enumerate}[label=\alph*)]
    \item \texttt{CREATE TABLE Aluno (matricula INT PRIMARY KEY, nome VARCHAR(100));}
    \item \texttt{CREATE Aluno TABLE matricula PRIMARY KEY, nome;}
    \item \texttt{INSERT TABLE Aluno (matricula INT PRIMARY KEY, nome VARCHAR(100));}
    \item \texttt{SELECT TABLE Aluno (matricula INT PRIMARY KEY, nome VARCHAR(100));}
\end{enumerate}

\item Considere a consulta:

\[
\texttt{SELECT nome FROM Aluno ORDER BY nome;}
\]

Ela retorna:

\begin{enumerate}[label=\alph*)]
    \item os nomes dos alunos ordenados por nome.
    \item os nomes dos alunos agrupados por curso.
    \item apenas alunos sem nome.
    \item a quantidade de alunos por nome.
\end{enumerate}

\item Considere a consulta:

\[
\texttt{SELECT COUNT(DISTINCT curso) FROM Aluno;}
\]

Ela retorna:

\begin{enumerate}[label=\alph*)]
    \item a quantidade de cursos distintos presentes na tabela Aluno.
    \item todos os alunos sem repetição.
    \item a soma das matrículas.
    \item a média de alunos por curso.
\end{enumerate}

\item Considere a tabela \texttt{Aluno}:
\[
\begin{array}{|c|c|c|}
\hline
\text{matricula} & \text{nome} & \text{curso}\\
\hline
1 & \text{Ana} & \text{BTI}\\
2 & \text{Bruno} & \text{Computacao}\\
3 & \text{Carla} & \text{BTI}\\
\hline
\end{array}
\]
Qual é o resultado da consulta abaixo?
\[
\texttt{SELECT nome FROM Aluno WHERE curso = 'BTI' ORDER BY nome;}
\]

\begin{enumerate}[label=\alph*)]
    \item Ana e Carla.
    \item Bruno apenas.
    \item BTI e Computacao.
    \item Nenhuma linha.
\end{enumerate}

\item Usando a tabela \texttt{Aluno} da questão anterior, qual é o resultado de:
\[
\texttt{SELECT COUNT(*) FROM Aluno WHERE curso = 'BTI';}
\]

\begin{enumerate}[label=\alph*)]
    \item 1.
    \item 2.
    \item 3.
    \item BTI.
\end{enumerate}

\item Considere as tabelas:
\[
Aluno(\text{matricula}, \text{nome}, \text{idCurso})
\]
\[
Curso(\text{idCurso}, \text{nomeCurso})
\]
com os dados:
\[
\begin{array}{|c|c|c|}
\hline
\text{matricula} & \text{nome} & \text{idCurso}\\
\hline
1 & \text{Ana} & 10\\
2 & \text{Bruno} & 20\\
\hline
\end{array}
\quad
\begin{array}{|c|c|}
\hline
\text{idCurso} & \text{nomeCurso}\\
\hline
10 & \text{BTI}\\
20 & \text{Computacao}\\
\hline
\end{array}
\]
Qual consulta retorna \texttt{Ana, BTI} e \texttt{Bruno, Computacao}?

\begin{enumerate}[label=\alph*)]
    \item \texttt{SELECT A.nome, C.nomeCurso FROM Aluno A, Curso C WHERE A.idCurso = C.idCurso;}
    \item \texttt{SELECT A.nome, C.nomeCurso FROM Aluno A, Curso C;}
    \item \texttt{SELECT nome, nomeCurso FROM Aluno WHERE idCurso = nomeCurso;}
    \item \texttt{DELETE FROM Curso WHERE idCurso = idCurso;}
\end{enumerate}

\item Considere a tabela \texttt{Matricula}:
\[
\begin{array}{|c|c|c|}
\hline
\text{idAluno} & \text{disciplina} & \text{nota}\\
\hline
1 & \text{BD} & 8.0\\
2 & \text{BD} & 7.0\\
3 & \text{PDS} & 9.0\\
\hline
\end{array}
\]
Qual consulta retorna apenas disciplinas com mais de um aluno matriculado?

\begin{enumerate}[label=\alph*)]
    \item \texttt{SELECT disciplina FROM Matricula WHERE COUNT(*) > 1;}
    \item \texttt{SELECT disciplina FROM Matricula GROUP BY COUNT(*) HAVING disciplina > 1;}
    \item \texttt{SELECT disciplina FROM Matricula GROUP BY disciplina HAVING COUNT(*) > 1;}
    \item \texttt{SELECT disciplina FROM Matricula HAVING nota > 1;}
\end{enumerate}

\item Considere:
\[
Aluno(\text{matricula}, \text{nome}, \text{idCurso})
\]
\[
Curso(\text{idCurso}, \text{nomeCurso}, \text{centro})
\]
Qual consulta retorna os nomes dos alunos de cursos do centro \texttt{'CT'}?

\begin{enumerate}[label=\alph*)]
    \item \texttt{SELECT nome FROM Aluno WHERE centro = 'CT';}
    \item \texttt{SELECT nome FROM Aluno WHERE idCurso IN (SELECT idCurso FROM Curso WHERE centro = 'CT');}
    \item \texttt{SELECT nomeCurso FROM Curso WHERE nome IN Aluno;}
    \item \texttt{SELECT centro FROM Aluno GROUP BY nome;}
\end{enumerate}

\item A consulta abaixo está incorreta em SQL padrão:
\[
\texttt{SELECT curso, COUNT(*) FROM Aluno WHERE COUNT(*) > 5 GROUP BY curso;}
\]
Qual é o principal erro?

\begin{enumerate}[label=\alph*)]
    \item Condições sobre agregações devem ser colocadas em \texttt{HAVING}, não em \texttt{WHERE}.
    \item \texttt{COUNT(*)} nunca pode ser usado com \texttt{GROUP BY}.
    \item \texttt{WHERE} só pode aparecer depois de \texttt{GROUP BY}.
    \item A consulta deveria usar \texttt{DELETE}.
\end{enumerate}

\item Se \texttt{Aluno} possui 3 linhas e \texttt{Curso} possui 2 linhas, a consulta abaixo retorna quantas combinações, assumindo ausência de filtro?
\[
\texttt{SELECT * FROM Aluno, Curso;}
\]

\begin{enumerate}[label=\alph*)]
    \item 2.
    \item 3.
    \item 5.
    \item 6.
\end{enumerate}

\item Considere o comando:
\[
\texttt{UPDATE Aluno SET curso = 'BTI';}
\]
Sem cláusula \texttt{WHERE}, o comando:

\begin{enumerate}[label=\alph*)]
    \item atualiza apenas a primeira tupla.
    \item gera sempre erro de sintaxe.
    \item atualiza todas as tuplas da tabela \texttt{Aluno}.
    \item remove todas as tuplas da tabela \texttt{Aluno}.
\end{enumerate}

\item Considere que \texttt{Aluno.idCurso} é chave estrangeira referenciando \texttt{Curso.idCurso}. Se tentarmos inserir um aluno com \texttt{idCurso = 99}, mas não existe curso 99, o que ocorre?

\begin{enumerate}[label=\alph*)]
    \item A inserção sempre é aceita.
    \item Há violação de integridade referencial.
    \item O SGBD cria automaticamente o curso 99.
    \item A chave estrangeira vira chave primária.
\end{enumerate}

\item Se \texttt{Aluno.matricula} é chave primária e já existe uma tupla com \texttt{matricula = 1}, inserir outra tupla com \texttt{matricula = 1} causa:

\begin{enumerate}[label=\alph*)]
    \item violação de chave primária.
    \item criação automática de nova tabela.
    \item ordenação dos dados.
    \item remoção da tupla antiga obrigatoriamente.
\end{enumerate}

\item Qual comando cria a tabela \texttt{Aluno} com chave estrangeira para \texttt{Curso}?

\begin{enumerate}[label=\alph*)]
    \item \texttt{CREATE TABLE Aluno (matricula INT PRIMARY KEY, nome VARCHAR(100), idCurso INT, FOREIGN KEY (idCurso) REFERENCES Curso(idCurso));}
    \item \texttt{CREATE TABLE Aluno FOREIGN KEY idCurso Curso;}
    \item \texttt{INSERT TABLE Aluno REFERENCES Curso(idCurso);}
    \item \texttt{SELECT FOREIGN KEY FROM Aluno TO Curso;}
\end{enumerate}

\item Se uma chave estrangeira de \texttt{Aluno} referencia \texttt{Curso} e não há ação de cascata configurada, tentar apagar um curso ainda referenciado por alunos normalmente:

\begin{enumerate}[label=\alph*)]
    \item apaga automaticamente todos os alunos em qualquer SGBD.
    \item é impedido por restrição de integridade referencial.
    \item transforma alunos em valores \texttt{NULL} obrigatoriamente.
    \item converte a tabela em visão.
\end{enumerate}

\item Considere:
\[
\texttt{SELECT nome FROM Aluno WHERE curso IS NULL;}
\]
Essa consulta retorna:

\begin{enumerate}[label=\alph*)]
    \item alunos cujo curso é diferente de \texttt{NULL}.
    \item todos os alunos, pois \texttt{NULL} é ignorado.
    \item alunos cujo curso está ausente ou desconhecido.
    \item erro, pois \texttt{IS NULL} não existe em SQL.
\end{enumerate}

\item Considere a tabela \texttt{Aluno}:
\[
\begin{array}{|c|c|}
\hline
\text{nome} & \text{nota}\\
\hline
Ana & 8.0\\
Bruno & 9.5\\
Carla & 7.0\\
\hline
\end{array}
\]
Qual aluno aparece primeiro no resultado?
\[
\texttt{SELECT nome FROM Aluno ORDER BY nota DESC;}
\]

\begin{enumerate}[label=\alph*)]
    \item Bruno.
    \item Ana.
    \item Carla.
    \item O resultado não pode ser ordenado por nota.
\end{enumerate}

\item Qual consulta retorna alunos que não possuem matrícula em nenhuma disciplina, considerando:
\[
Aluno(\text{matricula}, \text{nome})
\]
\[
Matricula(\text{matricula}, \text{disciplina})
\]

\begin{enumerate}[label=\alph*)]
    \item \texttt{SELECT nome FROM Aluno WHERE matricula NOT IN (SELECT matricula FROM Matricula);}
    \item \texttt{SELECT nome FROM Aluno WHERE COUNT(*) = 0;}
    \item \texttt{SELECT nome FROM Matricula WHERE disciplina IS NULL;}
    \item \texttt{DELETE FROM Aluno WHERE matricula IN Matricula;}
\end{enumerate}

\item Considere:
\[
\texttt{SELECT curso, AVG(nota) FROM Aluno GROUP BY curso;}
\]
Essa consulta retorna:

\begin{enumerate}[label=\alph*)]
    \item uma média de nota para cada curso.
    \item uma única média geral da tabela inteira, sem cursos.
    \item todas as notas individuais, sem agregação.
    \item uma nova tabela física chamada curso.
\end{enumerate}

\item Qual alternativa representa corretamente uma consulta com junção explícita?

\begin{enumerate}[label=\alph*)]
    \item \texttt{SELECT A.nome, C.nomeCurso FROM Aluno A JOIN Curso C ON A.idCurso = C.idCurso;}
    \item \texttt{SELECT A.nome JOIN C.nomeCurso FROM Aluno ON Curso;}
    \item \texttt{JOIN Aluno, Curso SELECT nome;}
    \item \texttt{SELECT nome FROM Aluno WHERE JOIN Curso;}
\end{enumerate}


\end{enumerate}

\newpage

\section*{Gabarito}

\begin{multicols}{2}
\begin{enumerate}[label=\textbf{Q\arabic*.}]
    \item a
    \item b
    \item b
    \item a
    \item b
    \item a
    \item a
    \item b
    \item c
    \item b
    \item a
    \item a
    \item b
    \item a
    \item a
    \item b
    \item a
    \item b
    \item b
    \item a
    \item a
    \item a
    \item a
    \item a
    \item b
    \item c
    \item a
    \item a
    \item b
    \item a
    \item b
    \item a
    \item a
    \item a
    \item a
    \item a
    \item a
    \item a
    \item a
    \item a
    \item a
    \item a
    \item a
    \item a
    \item a
    \item a
    \item a
    \item a
    \item a
    \item a
    \item a
    \item V, V, F, V
    \item V, V, V, F
    \item V, V, V, F
    \item V, V, V, F
    \item V, V, V, F
    \item V, V, V, F
    \item V, V, F, V
    \item V, V, V, F
    \item a
    \item a
    \item a
    \item a
    \item a
    \item a
    \item a
    \item a
    \item a
    \item a
    \item a
    \item b
    \item a
    \item c
    \item b
    \item a
    \item d
    \item c
    \item b
    \item a
    \item a
    \item b
    \item c
    \item a
    \item a
    \item a
\end{enumerate}
\end{multicols}

\end{document}
